% Adapted from David Levinson's official Findings template (CC BY 4.0).
% The unmodified source and retrieval hashes are in template_original/.
\documentclass{article}
\usepackage[english]{babel}
\usepackage[letterpaper,top=2cm,bottom=2cm,left=3cm,right=3cm,marginparwidth=1.75cm]{geometry}
\usepackage{amsmath,graphicx,booktabs,array}
\usepackage[sort&compress,round,semicolon,authoryear]{natbib}
\usepackage[figurename=Fig.,labelfont=bf,labelsep=period]{caption}
\usepackage[colorlinks=true,allcolors=blue]{hyperref}
\usepackage[none]{hyphenat}
\emergencystretch=4em
\graphicspath{{../figures/}}
\newcommand{\pending}[1]{\textbf{[RESULT PENDING: #1]}}
\newcommand{\studyfigure}[2]{\IfFileExists{../figures/#1.pdf}{\includegraphics[width=\linewidth]{#1.pdf}}{\fbox{\parbox[c][4cm][c]{0.92\linewidth}{\centering Figure awaiting generated results: #2}}}}
\title{Can LLM Personas Reproduce Ground-Mode Co-Use on Air Journeys? A Public-Data Benchmark for Germany}
\author{}
\date{November 24, 2026}
\hypersetup{pdftitle={Can LLM Personas Reproduce Ground-Mode Co-Use on Air Journeys?},pdfsubject={Aggregate MiD 2023 benchmark},pdfkeywords={large language models,personas,air travel,national travel survey,mode co-use}}

\begin{document}
\maketitle
\begin{abstract}
We adapt LLM-PDM to estimate ground-mode co-use on overnight air journeys using Germany's MiD 2023. Public tables define 23 single-attribute cohort personas; four fully evaluated models predict car, local-public-transport, and long-distance-rail probabilities. Overall-rate guidance reduced error for every model, but only guided GPT-4.1 outperformed an observed-rate constant baseline: mean absolute error 2.93 versus 4.16 percentage points. Unguided errors were 9.44--25.08 points. These results demonstrate model-dependent aggregate agreement, not individual simulation or airport-access validity.
\end{abstract}
\textit{Keywords:} large language models; personas; air journeys; national travel survey; multimodality; benchmarking

\section{Questions}
Persona-based LLM surveys can reproduce selected mobility distributions \citep{tzachristas2026}, but agreement with overall averages does not establish fidelity within groups \citep{bisbee2024}. Travel-choice applications also evaluate individual predictions using labeled records \citep{mo2026,liu2026}. We ask a narrower question: can LLMs recover publicly observable differences in ground-mode co-use among air journeys, beyond simply repeating overall rates? Does supplying those rates improve agreement? Our contribution is a reproducible, aggregate-only adaptation of LLM-PDM that separates unconditional guidance from withheld conditional targets (Figure~\ref{fig:framework}).

\section{Methods}
We used public Mobilit\"at in Tabellen cross-tabulations from MiD 2023 \citep{mit2023}. The overnight-journey module covers self-responding residents of Germany aged at least 14, reporting up to three journeys during the preceding three months; fieldwork occurred in 2023--2024 \citep{midmethods2025,midhandbook2025}. We selected the recorded main-mode category aircraft. Each observation is a journey, not a unique passenger. The multiple-response question covers outward and return travel together. Car, local bus/rail, and long-distance train are three nonexclusive outcomes; none identifies a particular airport-access leg.

Before study predictions, we fixed a screening rule for 14 candidate dimensions. Substantive strata required at least 100 unweighted journeys; dimensions required at least two eligible strata covering 80\% of weighted aircraft journeys. We ranked mean normalized Bernoulli information gain across the three outcomes and selected two traveler and two journey dimensions. All eligible strata became separate cohort personas. Missing categories were excluded, with coverage retained. Earlier exploratory age tables had been inspected; this was not an externally registered preregistration.

Unlike LLM-PDM's joint-persona sampling, each prompt described one attribute only and requested three probabilities directly. No joint attributes, synthetic individuals, fine-tuning, or iterative target matching were used. Gemma~3~4B, Gemma~3~12B, GPT-4.1~mini, Qwen3~235B~A22B, and GPT-4.1 were queried through OpenRouter in three repeats at temperature~0.7. Both arms withheld cohort rates; one supplied overall rates. No-persona controls retained the population and task description. Model tiers are not a controlled parameter-size comparison.

Primary error was mean absolute error (MAE) of the three-repeat mean probabilities, in percentage points, weighted by stratum prevalence within each dimension and equally across dimensions and outcomes. Supplementary checks include RMSE, Bernoulli Jensen--Shannon divergence, and repeat variability. The observed overall-rate baseline uses no personas. Screening and evaluation share survey aggregates: this is an in-sample descriptive benchmark, without survey-design confidence intervals or individual-level validation.

\begin{figure}[htbp]
\centering
\studyfigure{01_framework}{Study framework}
\caption{Aggregate adaptation of LLM-PDM. Conditional observations select cohorts and score predictions; only overall rates enter the anchored prompts.}
\label{fig:framework}
\end{figure}

\section{Findings}
The aircraft cohort contained 9,475 reported overnight journeys. Overall car, local-public-transport, and long-distance-rail co-use was 37.7\%, 21.6\%, and 12.4\%. Household car count, residential spatial type, household companions, and nights away yielded 23 profiles covering 92.5--100\% of weighted journeys within each dimension (Figure~\ref{fig:selection}). Car co-use was 15.0\% in households without cars and 58.4\% in households with at least four cars; corresponding local-public-transport rates were 40.1\% and 9.3\%.

Four models completed all 144 calls each. Gemma~3~12B returned only 11 valid outputs, with 98 rate-limited failures and 35 invalid schemas; it was excluded before comparative scoring. Overall-rate guidance lowered every complete model's MAE: unanchored errors ranged from 9.44 to 25.08 points, versus 2.93--6.80 when anchored (Figure~\ref{fig:results}). Only anchored GPT-4.1 personas beat the 4.16-point constant baseline, reaching 2.93 points; its no-persona control scored 4.10. Persona conditioning worsened anchored predictions for the other three models. GPT-4.1's individual-repeat MAEs were 2.89--3.26 points. Its largest remaining error concerned long-distance rail in small-town/village areas within urban regions: 8.3\% predicted versus 23.3\% observed. Generated rationales are hypotheses, not evidence of respondents' motives.

Overall-rate guidance helped, but persona-specific improvement beyond simple rates was model-dependent. Agreement supports only the tested aggregate margins. Unknown pretraining exposure, reused selection targets, unavailable survey-design covariance, and unobserved joint characteristics limit interpretation. These data cannot validate individual behavior, causal mechanisms, Munich passengers, or airport-access mode choice. They provide an openly auditable test before such stronger applications are attempted.

\begin{figure}[htbp]
\centering
\studyfigure{02_persona_selection}{Selected cohort profiles}
\caption{Fixed-rule dimension selection and observed cohort co-use. Strata are separate marginal profiles, not combinations of attributes.}
\label{fig:selection}
\end{figure}

\begin{figure}[htbp]
\centering
\studyfigure{03_model_results}{Model performance}
\caption{A: complete-model MAE versus the overall-rate baseline. Points score mean probabilities; bars show single-repeat MAE ranges, not confidence intervals. B--C: GPT-4.1's 69 cohort/mode predictions before and after guidance. Gemma~3~12B was operationally incomplete.}
\label{fig:results}
\end{figure}

\clearpage
\section*{Acknowledgments and availability}
OpenAI Codex assisted with code and manuscript preparation. The evaluated models and their research role are described above. Public source responses, prompts, outputs, and analysis code are provided in the \href{https://github.com/jtzach/air-passenger-LLM-PDM-study}{replication repository}.
\bibliographystyle{chicago}
\bibliography{refs}
\end{document}
